\documentclass[10pt,a4paper]{report}
\usepackage[utf8]{inputenc}
\usepackage[spanish]{babel}
\usepackage{amsmath}
\usepackage{amsfonts}
\usepackage{amssymb}
\usepackage{steinmetz}
\usepackage{circuitikz}
\usepackage{halloweenmath}
\usepackage[normalem]{ulem}
%\usepackage{steinmetz}
%\usepackage[american,siunitx]{circuitikz}
\usetikzlibrary{arrows,shapes,calc,positioning}
%%%%%%%%%%%%%%%%%%%%%%%%%%%%%%%%%%%%%%%%%%%%%%%%
\usepackage{listings}
\usepackage{xcolor}

\definecolor{codegreen}{rgb}{0,0.6,0}
\definecolor{codegray}{rgb}{0.5,0.5,0.5}
\definecolor{codepurple}{rgb}{0.58,0,0.82}
\definecolor{backcolour}{rgb}{0.95,0.95,0.92}

\lstdefinestyle{mystyle}{
	backgroundcolor=\color{backcolour},   
	commentstyle=\color{codegreen},
	keywordstyle=\color{magenta},
	numberstyle=\tiny\color{codegray},
	stringstyle=\color{codepurple},
	basicstyle=\ttfamily\footnotesize,
	breakatwhitespace=false,         
	breaklines=true,                 
	captionpos=b,                    
	keepspaces=true,                 
	numbers=left,                    
	numbersep=5pt,                  
	showspaces=false,                
	showstringspaces=false,
	showtabs=false,                  
	tabsize=2
}
\lstset{style=mystyle}
%%%%%%%%%%%%%%%%%%%%%%%%%%%%%%%%%%%%%%%%%%%%%%%%
\usepackage{graphicx}
\usepackage{multirow}
\usepackage[hyperindex=true,breaklinks=true,colorlinks=true,linkcolor=blue]{hyperref}
\usepackage{multimedia}
\usepackage{subcaption}
\addto\captionsspanish{
\renewcommand{\tablename}{Tabla}
\renewcommand{\listtablename}{\'Indice de Tablas}
}

\usepackage{makeidx}
\usepackage[left=2cm,right=2cm,top=2cm,bottom=2cm]{geometry}
\begin{document}
\section{sofware necesario}

\begin{enumerate}
	\item La versión de Ubuntu que tenemos es la Ubuntu 24.04.4 LTS
	\begin{verbatim}
	https://ubuntu.com/download/desktop
	\end{verbatim}
	Es la última LTS, así que creo que el la que debemos instalar
	\item Visualstudio Code.
	\begin{verbatim}
	https://code.visualstudio.com/thank-you?dv=linux64_deb
	\end{verbatim}
Es un paquete deb, para instalarlo ir a la línea de comandos:
\begin{verbatim}
sudo apt install ./Descargas/code_*.deb	
\end{verbatim}

	
	\item STM32CubeCLT.
	
	https://www.st.com/en/development-tools/stm32cubeclt.html

	Tiene todas las herramientas de ST, necesarias para compilar, enlazar etc.
	Descargar el que pone Debian Linux installer
	Descomprimir
	\begin{verbatim}
		unzip st-stm32cubeclt_1.21.0_27995_20260219_1804_amd64.deb_bundle.sh.zip
	\end{verbatim}
	Dar Permisos
	\begin{verbatim}
		chmod +x st-stm32cubeclt_1.21.0_27995_20260219_1804_amd64.deb_bundle.sh
	\end{verbatim}
Lanzar el instalador
\begin{verbatim}
sudo ./st-stm32cubeclt_1.21.0_27995_20260219_1804_amd64.deb_bundle.sh	
\end{verbatim}
Ten sacará en el terminal la licencia de uso, ir hasta el final y aceptarla.
Debería instalarlo en:
\begin{verbatim}
	/opt/st/stm32cubeclt_1.21.0/
\end{verbatim}

\item stm32cubemx. Realmente no lo tendrían que usar, es una herramienta para configurar la placa pero en algún momento puede ser de gran ayuda.
\begin{verbatim}
https://www.st.com/en/development-tools/stm32cubemx.html
\end{verbatim}

El resto de las herramientas, se pueden instalar directamente a través de Visualstudio Code Mediante las extensiones para SMT32,(barra lateral izquierda, icono de los cuadrados). Para ello hay que abrir VSCode y buscar las extensiones de STM32CUBE. El símbolo es una mariposa. Es importante emplear las oficiales. 

\item Python y los módulos de Python: serial, struct, csv, sys, time, collections, datetime, pyqtgraph.
\item (Opcional) Si se va a hacer el análisis de los datos y los modelos del motor en python, entonces es conveniente instalar también los módulos: pandas, numpy, matplotlib.pyplot y control

Nota: creo que lo mas fácil es instalar anaconda.
\begin{verbatim}
https://www.anaconda.com/download
\end{verbatim}

Varios de los módulos de Python necesarios vienen instalados por defecto con Anaconda y el resto son fácilmente instalables empleando \texttt{conda}.

\item Descargar el proyecto de github
\begin{verbatim}
git clone https://github.com/juanjimenez/st_projects.git
\end{verbatim}
La carpeta que hay que abrir en VS Code es la carpeta motor. 
De momento, probamos a ver como va,y luego ya la podaremos para dejar a los alumnos solo lo que de verdad necesitan.
\end{enumerate}

\end{document}